% source: 0d34333946fc2a5557afa862464cb9b87b7400d1
\chapter{SLH-DSA（SPHINCS\textsuperscript{+}，FIPS 205）：算法详解}
\chaptermark{SLH-DSA / SPHINCS\textsuperscript{+}}
\label{ch:slhdsa}

\section{什么是 SPHINCS\textsuperscript{+}？}

\textbf{SPHINCS\textsuperscript{+}} 是 NIST 在 FIPS 205~\cite{fips205,sphincsplus2019} 中标准化为 \textbf{SLH-DSA} 的基于哈希的签名方案，它是所有后量子标准中最保守的一个。上一部分的格签名依赖于格问题被假定的困难性，而 SPHINCS\textsuperscript{+} \emph{仅仅依赖于哈希函数的安全性}——这是密码学中被理解得最透彻的单一假设。即使明天所有的格假设和编码假设都被攻破，一个安全的哈希函数仍然能够产生安全的签名；这正是该方案族被保留在标准组合中作为对冲的原因。

该方案是一条简短谱系的终点：无状态设计 \emph{SPHINCS} 于 2015 年问世，在 2017--2019 年间为参加 NIST 竞赛而被强化为 \emph{SPHINCS\textsuperscript{+}}~\cite{sphincsplus2019}，并于 2024 年在 FIPS 205 中最终定型为 \emph{SLH-DSA}。这三个名称指代的基本上是同一个构造；本书用 SLH-DSA 指称标准化后的算法，用 SPHINCS\textsuperscript{+} 指称其前身方案。用一句话概括，它是\emph{一棵由一次性与少次哈希签名构成的巨型超树，其排布方式使得永远无需保存任何签名状态}。

\paragraph{为什么“无状态”是难点。}
SLH-DSA 中的 “SL” 代表\emph{无状态}（stateless），而这正是该设计竭力争取的性质。较早的基于哈希的方案，如 XMSS 和 LMS，是\emph{有状态}的：每个一次性密钥只能恰好使用一次，因此签名者必须记住哪些叶子已被用掉，而一次疏忽——恢复了一个虚拟机快照、复制了一个密钥文件——就会重用一次性密钥，并可能造成灾难性后果。SPHINCS\textsuperscript{+} 彻底消除了状态。它使用消息随机化和一个极其庞大的超树地址空间，使意外的地址碰撞可以忽略。这并不意味着允许重用地址：WOTS 叶子必须保持唯一，而 FORS 的安全性会随有限的签名次数以及任何意外碰撞而退化。无状态的代价是尺寸：SPHINCS\textsuperscript{+} 的签名是所有标准中最大的，每个签名有数 KB。该方案可以用 \textbf{SHA-2} 或 \textbf{SHAKE} 实例化，有 \emph{simple}（简单）和 \emph{robust}（稳健）两种模式，并分为下文参数一节中详述的、在尺寸与速度之间取舍的 \textbf{s} 族和 \textbf{f} 族。

\section{学习目标}

前两章从零开始构建了基于哈希的签名的各个组件：一次性签名与少次签名（第~\ref{ch:hashots}~章），以及使它们成为无状态、可多次使用的 Merkle 树与超树机制（第~\ref{ch:hashmerkle}~章）。本章完全按照标准的规定把这些组件组装起来。读完本章后，你应当能够：

\begin{itemize}
  \item 解释为何 SPHINCS\textsuperscript{+} 是最保守的标准，以及“无状态”带来了什么、付出了什么代价；
  \item 回忆 WOTS\textsuperscript{+}、FORS、Merkle 树和超树是如何组合在一起的；
  \item 解释每个秘密值如何由少数几个种子和一个地址确定性地派生出来；
  \item 理解 SLH-DSA 的密钥生成、签名和验证算法；
  \item 读懂各参数族及其在尺寸与速度之间的权衡。
\end{itemize}

\section{回顾：组装起来的各个组件}
值得在一处重新陈述这一阶梯，因为 SLH-DSA 恰恰就是它的组装：

\begin{itemize}
  \item \textbf{WOTS\textsuperscript{+}}（第~\ref{ch:hashots}~章）用哈希链和校验和对一个值签名；
  \item 一棵 \textbf{Merkle 树}（第~\ref{ch:hashmerkle}~章）把 $2^h$ 个 WOTS\textsuperscript{+} 密钥绑定在一个根之下，并由一条路径进行认证——一棵这样的树就是一个 XMSS 实例；
  \item 一棵\textbf{超树}把 $d$ 层这样的树堆叠起来，使每个根都由上一层签名，只把顶端的根作为公钥公开；
  \item \textbf{FORS} 作为一种\emph{少次}签名对实际的消息摘要签名，而它的公钥正是超树最底层所签名的对象。
\end{itemize}

\begin{figure}[H]
\centering
\begin{tikzpicture}[
  scale=0.82,
  mnode/.style={circle, draw=pqcgray, fill=white, minimum size=3.2mm, inner sep=0pt},
  wnode/.style={rectangle, draw=pqcgray, fill=white, minimum size=2.8mm, inner sep=0pt},
  wused/.style={rectangle, draw=pqcgray, fill=black, minimum size=2.8mm, inner sep=0pt},
  fnode/.style={regular polygon, regular polygon sides=6, draw=pqcgray, fill=pqcgray!40, minimum size=6mm, inner sep=0pt},
  sig/.style={pqcgray, densely dotted, thick},
  font=\scriptsize
]
  \foreach \i/\yb/\lab in {0/2.2/{layer $0$}, 1/4.2/{layer $1$}, 2/6.2/{layer $d-1=2$}} {
    \fill[pqclight] (-2,\yb-0.18) -- (2,\yb-0.18) -- (0,\yb+1.62) -- cycle;
    \draw[pqcgray!55] (-2,\yb-0.18) -- (2,\yb-0.18) -- (0,\yb+1.62) -- cycle;
    \node[pqcgray, anchor=west] at (2.15,\yb+0.95) {\lab};
    \node[mnode] (r\i) at (0,\yb+1.4) {};
    \node[mnode] (mL\i) at (-1,\yb+0.7) {};
    \node[mnode] (mR\i) at (1,\yb+0.7) {};
    \node[wnode] (a\i) at (-1.5,\yb) {};
    \node[wused] (leaf\i) at (-0.5,\yb) {};
    \node[wnode] (c\i) at (0.5,\yb) {};
    \node[wnode] (d\i) at (1.5,\yb) {};
    \draw[pqcgray] (r\i) -- (mL\i) (r\i) -- (mR\i)
                   (mL\i) -- (a\i) (mL\i) -- (leaf\i)
                   (mR\i) -- (c\i) (mR\i) -- (d\i);
    \draw[pqcorange, very thick] (leaf\i) -- (mL\i) -- (r\i);
  }
  \node[anchor=west, font=\bfseries] at ($(r2)+(0.12,0)$) {$PK.\mathrm{root}$};
  \draw[sig] (r0) -- (leaf1) node[midway, right=0.8mm] {WOTS\ensuremath{^{+}} sig.};
  \draw[sig] (r1) -- (leaf2) node[midway, right=0.8mm] {WOTS\ensuremath{^{+}} sig.};
  \node[fnode] (fors) at (-0.5,1.1) {};
  \node[anchor=west] at ($(fors)+(0.45,0)$) {FORS public key};
  \node[ellipse, draw=pqcgray, fill=pqclight, minimum width=2.1cm, minimum height=0.72cm] (msg) at (-0.5,-0.2) {Message};
  \draw[sig] (fors) -- (leaf0) node[midway, right=0.8mm] {WOTS\ensuremath{^{+}} sig.};
  \draw[sig] (msg) -- (fors) node[midway, right=0.8mm] {FORS sig.};
  \draw[pqcgray, rounded corners] (3.7,3.35) rectangle (7.7,5.45);
  \node[mnode] at (4.05,4.95) {}; \node[anchor=west] at (4.35,4.95) {Merkle tree node};
  \node[wnode] at (4.05,4.40) {}; \node[anchor=west] at (4.35,4.40) {WOTS\ensuremath{^{+}} public key};
  \node[fnode, minimum size=4.5mm] at (4.05,3.85) {}; \node[anchor=west] at (4.35,3.85) {FORS public key};
\end{tikzpicture}
\caption{一个 SLH-DSA 签名：$d=3$ 层的完整超树。消息由 \textbf{FORS}（一种少次签名）签名；FORS 公钥由第 $0$ 层 Merkle 树的一个 WOTS\textsuperscript{+} 叶子签名；每棵树的根由上一层树的一个 WOTS\textsuperscript{+} 叶子签名；只有顶端的根 $PK.\mathrm{root}$ 被公布。每棵树中被使用的叶子（实心）及其通往根的路径被突出显示——这正是把叶子与根联系起来的认证数据。仿照 FIPS 205~\cite{fips205} 的风格重绘。}
\label{fig:slhdsa-signature}
\end{figure}

图~\ref{fig:slhdsa-signature} 画出了一个完整的 SLH-DSA 签名：从底部的消息一直到已公布的根的整条一次性签名与少次签名链（同一超树的模块级视图见图~\ref{fig:hypertree}）。剩下要做的是明确\emph{每个秘密值从何而来}，并写出三个算法。

\section{真正的输入：种子与地址}
与 ML-KEM 和 ML-DSA 一样，SLH-DSA 一旦固定了种子便是确定性的。一切都由三个 $n$ 字节的值派生而来：

\begin{itemize}
  \item \texttt{SK.seed}——秘密；每个 WOTS\textsuperscript{+} 和 FORS 秘密值都通过一个 PRF 由它生成；
  \item \texttt{SK.prf}——秘密；一个 PRF 密钥，用于派生每个签名各自的随机化值 $R$；
  \item \texttt{PK.seed}——公开；对每次哈希进行调整，使结构上相同的节点互不相同。
\end{itemize}

使众多哈希保持彼此区分的机制是\emph{地址}（\texttt{ADRS}）：一个结构化的标签，准确标识某次哈希属于哪棵树、哪一层、哪个叶子以及链中的哪个位置。每次哈希调用都接收公开种子和一个地址，因此在不同位置上结构相同的两次计算绝不会产生相同的值。这正是挫败多目标攻击的关键，也是为什么 \texttt{PK.seed} 虽然公开，却不可或缺。

\section{密钥生成}
密钥生成采样三个种子，并计算概括整棵超树的那一个值：其顶层 Merkle 树的根。

\begin{algorithm}[H]
\caption{SLH-DSA.KeyGen（教学形式）}
\label{alg:slhdsa-keygen}
\begin{algorithmic}[1]
\Require 用于三个 $n$ 字节种子的随机性
\Ensure 公钥 $(\texttt{PK.seed},\texttt{PK.root})$；秘密钥还存储 $(\texttt{SK.seed},\texttt{SK.prf})$
\State 均匀采样 \texttt{SK.seed}、\texttt{SK.prf}、\texttt{PK.seed}
\State $\texttt{PK.root} \gets$ 顶层 Merkle 树的根，由 \texttt{SK.seed} 和 \texttt{PK.seed} 通过带地址的哈希派生
\State \Return 公钥 $(\texttt{PK.seed},\texttt{PK.root})$，秘密钥 $(\texttt{SK.seed},\texttt{SK.prf},\texttt{PK.seed},\texttt{PK.root})$
\end{algorithmic}
\end{algorithm}

\section{签名与验证}
签名过程把消息（连同一个随机化值）哈希为摘要，将摘要拆分为一条 FORS 消息以及一对用于选择底层树和叶子的索引，用 FORS 对摘要签名，然后沿超树向上对 FORS 公钥签名。

\begin{algorithm}[H]
\caption{SLH-DSA.Sign（教学形式）}
\label{alg:slhdsa-sign}
\begin{algorithmic}[1]
\Require 秘密钥，格式化后的消息 $M'$，可选的 $n$ 字节随机化值 $opt\_rand$
\Ensure 签名 $\sigma = (R,\ \sigma_{\mathrm{FORS}},\ \sigma_{\mathrm{HT}})$
\State $R \gets \textsc{PRF\_msg}(\texttt{SK.prf}, opt\_rand, M')$ \Comment{$opt\_rand$ 为零时是确定性的；为随机输入时是对冲（hedged）的}
\State $digest \gets H(R \,\|\, \texttt{PK.seed} \,\|\, \texttt{PK.root} \,\|\, M')$
\State 将 $digest$ 拆分为 FORS 消息 $md$ 与索引 $(idx_{\mathrm{tree}}, idx_{\mathrm{leaf}})$
\State $\sigma_{\mathrm{FORS}} \gets \textsc{ForsSign}(md, \dots)$;\quad $pk_{\mathrm{FORS}} \gets$ 其公钥
\State $\sigma_{\mathrm{HT}} \gets \textsc{HyperTreeSign}(pk_{\mathrm{FORS}}, idx_{\mathrm{tree}}, idx_{\mathrm{leaf}}, \dots)$
\State \Return $(R,\ \sigma_{\mathrm{FORS}},\ \sigma_{\mathrm{HT}})$
\end{algorithmic}
\end{algorithm}

\begin{algorithm}[H]
\caption{SLH-DSA.Verify（教学形式）}
\label{alg:slhdsa-verify}
\begin{algorithmic}[1]
\Require 公钥 $(\texttt{PK.seed},\texttt{PK.root})$，消息 $M$，签名 $(R,\sigma_{\mathrm{FORS}},\sigma_{\mathrm{HT}})$
\Ensure \emph{接受}或\emph{拒绝}
\State $digest \gets H(R \,\|\, \texttt{PK.seed} \,\|\, \texttt{PK.root} \,\|\, M)$；拆分为 $md,(idx_{\mathrm{tree}},idx_{\mathrm{leaf}})$
\State $pk_{\mathrm{FORS}} \gets \textsc{ForsPkFromSig}(md, \sigma_{\mathrm{FORS}}, \dots)$
\State $root' \gets \textsc{HyperTreeRootFromSig}(pk_{\mathrm{FORS}}, \sigma_{\mathrm{HT}}, idx_{\mathrm{tree}}, idx_{\mathrm{leaf}}, \dots)$
\State \Return 若 $root' = \texttt{PK.root}$ 则\emph{接受}，否则\emph{拒绝}
\end{algorithmic}
\end{algorithm}

验证不过是一连串的哈希重新计算：由消息重建 FORS 公钥，利用所提供的认证路径沿超树向上折叠，只有当最终值等于已公布的根时才接受。这里没有需要保护的算术运算，也没有依赖于秘密的分支，这使得 SLH-DSA 异常容易以常数时间实现（第~\ref{ch:sidechannel}~章）。

\section{参数集}
SLH-DSA 在每个安全级别上都有两种版本：\textbf{s}（\emph{小}签名，签名较慢）和 \textbf{f}（\emph{快速}签名，签名较大），可用 SHA2 或 SHAKE 实例化。

\begin{table}[H]
\centering
\begin{tabular}{lccc}
\toprule
参数族 & $n$（字节） & NIST 级别 & 权衡 \\
\midrule
SLH-DSA-128\{s,f\} & 16 & 1 & 密钥最小 \\
SLH-DSA-192\{s,f\} & 24 & 3 & 均衡 \\
SLH-DSA-256\{s,f\} & 32 & 5 & 安全性最高 \\
\bottomrule
\end{tabular}
\caption{SLH-DSA 参数族。在每个级别内，\textbf{s} 变体使签名尺寸最小，\textbf{f} 变体使签名时间最短。}
\end{table}

签名很大——大约 $8$ 到 $50$ KB——这是仅依赖哈希所付出的代价。作为回报，其安全性论证是所有后量子方案中最简单、最保守的。

\section{SageMath 实验：带地址的 Merkle 认证}

这个仅使用哈希的小例子检验 SLH-DSA 验证中把认证路径折叠回公开根的那一部分。它使用八个叶子和一个基于 SHA-256 的玩具版哈希；它不是 FORS、WOTS\textsuperscript{+} 或完整 SLH-DSA 的实现。

\begin{lstlisting}[language=Python]
from hashlib import sha256
from operator import xor

n = 16
pk_seed = b"public seed for a toy tree"

def H(label, address, data=b""):
    return sha256(pk_seed + label + address + data).digest()[:n]

def leaf(index):
    return H(b"leaf", int(index).to_bytes(4, "big"))

def parent(left, right, level, index):
    address = (int(level).to_bytes(1, "big") +
               int(index).to_bytes(4, "big"))
    return H(b"node", address, left + right)

levels = [[leaf(i) for i in range(8)]]
for level in range(3):
    previous = levels[-1]
    levels.append([parent(previous[2*i], previous[2*i + 1], level, i)
                   for i in range(len(previous)//2)])
root = levels[-1][0]

leaf_index = 5
auth = [levels[level][xor(int(leaf_index >> level), 1)]
        for level in range(3)]

def root_from_path(node, index, path):
    for level, sibling in enumerate(path):
        parent_index = index >> 1
        node = (parent(node, sibling, level, parent_index)
                if index % 2 == 0 else
                parent(sibling, node, level, parent_index))
        index >>= 1
    return node

assert root_from_path(leaf(leaf_index), leaf_index, auth) == root
bad_auth = list(auth)
bad_auth[0] = bytes([xor(bad_auth[0][0], 1)]) + bad_auth[0][1:]
assert root_from_path(leaf(leaf_index), leaf_index, bad_auth) != root
print("addressed Merkle path accepts intact data and rejects a tampered path")
\end{lstlisting}

地址被包含在每次哈希调用之中，因此位于不同层或不同节点索引处的相同字节串会被域分离。这就是贯穿整个 SLH-DSA 的地址规范的小规模版本。

\section{一个完整的 SageMath 实现}
\label{sec:fips205-impl}

上面的 Merkle 实验在一个要紧的方面是玩具版的：它使用了一个虚构的地址格式。配套文件 \sagefile{fips205\_slhdsa.sage} 实现了真正的方案——FIPS~205 中全部二十五个编号算法，每个算法都是一个独立的函数，并标注了其算法编号，覆盖全部十二个获批的参数集。这意味着两个哈希族（SHA2 和 SHAKE）、全部三个安全级别、\textbf{s} 和 \textbf{f} 两种变体都包括在内，此外还有 WOTS\textsuperscript{+} 链、XMSS 树、超树、FORS、纯消息（pure）与预哈希（pre-hash）两种顶层算法，以及 32 字节的地址结构及其 22 字节的压缩形式。

\playground{\texttt{f} 参数集的端到端运行，以及 \texttt{s} 参数集（其签名需要数分钟）对照 NIST ACVP 签名的验证}

地址是一个小类，携带标准表 1 中的各个成员函数，因此算法 9 的递归树遍历读起来与规范完全一致：

\begin{lstlisting}[language=Python]
def xmss_node(self, skseed, i, z, pkseed, adrs):
    """Algorithm 9. Root of a Merkle subtree of WOTS+ public keys."""
    if z == 0:
        adrs.setTypeAndClear(WOTS_HASH)
        adrs.setKeyPairAddress(i)
        return self.wots_pkGen(skseed, pkseed, adrs)
    lnode = self.xmss_node(skseed, 2*i,     z - 1, pkseed, adrs)
    rnode = self.xmss_node(skseed, 2*i + 1, z - 1, pkseed, adrs)
    adrs.setTypeAndClear(TREE)
    adrs.setTreeHeight(z)
    adrs.setTreeIndex(i)
    return self._H(pkseed, adrs, lnode + rnode)
\end{lstlisting}

超树于是不过是把这种树堆叠 $d$ 层，每一层对其下一层的根签名：

\begin{lstlisting}[language=Python]
def ht_sign(self, M, skseed, pkseed, idxtree, idxleaf):
    """Algorithm 12. A hypertree signature: d chained XMSS signatures."""
    adrs = ADRS()
    adrs.setTreeAddress(idxtree)
    SIGtmp = self.xmss_sign(M, skseed, idxleaf, pkseed, adrs)
    SIG_HT = SIGtmp
    root   = self.xmss_pkFromSig(idxleaf, SIGtmp, M, pkseed, adrs)
    for j in range(1, self.d):
        idxleaf = idxtree % (1 << self.hp)   # low h' bits
        idxtree = idxtree >> self.hp
        adrs.setLayerAddress(j)
        adrs.setTreeAddress(idxtree)
        SIGtmp  = self.xmss_sign(root, skseed, idxleaf, pkseed, adrs)
        SIG_HT += SIGtmp
        if j < self.d - 1:
            root = self.xmss_pkFromSig(idxleaf, SIGtmp, root, pkseed, adrs)
    return SIG_HT
\end{lstlisting}

\paragraph{对于一个没有代数结构的方案，Sage 能贡献什么。}
SLH-DSA 完全由哈希函数构建，因此与 ML-KEM 和 ML-DSA 不同，这里没有可供计算机代数系统用来核对代码的环。Sage 所提供的则是关于其组合结构的精确整数推理。位于 \sagefile{fips205\_slhdsa.sage} 末尾的函数 \texttt{verify\_parameter\_consistency()} 针对全部十二个参数集，从定义方程重新推导出表 2：验证 $h = d\cdot h'$，从而超树确实有 $2^{h}$ 个叶子；验证由算法 1 得出的校验和长度 $len_2$ 与闭式 $\lfloor \log_w(len_1(w-1)) \rfloor + 1$ 一致，后者用 Sage 的精确整数对数计算；验证摘要长度 $m$ 足以提供 FORS 消息以及两个树索引；并验证已公布的密钥尺寸和签名尺寸分别等于 $2n$ 与 $(1 + k(1+a) + h + d\cdot len)\,n$。参数表恰恰是那种容易抄错、又难以察觉的东西，而这一函数把它变成了一项测试。

\paragraph{验证。}
该实现针对 NIST 的 FIPS~205 ACVP 向量~\cite{acvp} 进行了检验——包括密钥生成、签名生成和签名验证，覆盖全部十二个参数集的纯消息形式与预哈希形式——并逐字节地复现了这些向量。

有一个实际的注意事项是方案本身的性质，而非本实现的性质。使用 \textbf{s} 参数集签名需要重建整棵 XMSS 树：SLH-DSA-128s 每个签名需要遍历大约一百万次量级的哈希调用，在 Sage 中约需十六秒，而 256 比特的变体约需其两倍。因此，快速测试只对 \textbf{f} 参数集覆盖签名\emph{生成}，并报告它略去了多少组；\texttt{make kat-full} 则把它们包括在内。验证对每个参数集来说开销都很小，因此总是被完整检验。正是这种不对称性给两种变体带来了各自的名称。

\section{常见陷阱}

\paragraph{陷阱 1：重用一次性（WOTS\textsuperscript{+}）叶子。}
用同一条 WOTS\textsuperscript{+} 链对两条不同的消息签名，会揭示足以进行伪造的链上点（第~\ref{ch:hashots}~章）。地址机制和消息随机化使意外碰撞极不可能发生；WOTS 地址和 FORS 地址都不得被有意重用。

\paragraph{陷阱 2：以为 SLH-DSA 需要状态。}
较早的基于哈希的方案（XMSS、LMS）是\emph{有状态}的：签名者必须记住哪些叶子已被用掉，一次重用就是灾难性的。SLH-DSA 是\emph{无状态}的——由消息哈希来选择路径——这正是尽管签名更大，它仍被选中进行标准化的原因。

\paragraph{陷阱 3：忽视公开种子的作用。}
\texttt{PK.seed} 通过地址机制调整每一次哈希调用。没有它，不同树中结构相同的节点会得到相同的哈希值，从而使多目标攻击成为可能。它是公开的，但不可或缺。

\section{小结}
SLH-DSA 是把基于哈希的阶梯组装成单个无状态标准的结果：

\begin{itemize}
  \item 每个秘密值都由 \texttt{SK.seed} 和一个地址确定性地派生；\texttt{PK.seed} 使所有哈希保持互不相同；
  \item 一个签名是对消息摘要的 FORS 签名，经由一棵由 Merkle 树（XMSS）实例构成的超树向上传递，直到唯一公布的根；
  \item 验证纯粹是哈希运算，因此该方案在假设上保守，并且易于实现为常数时间。
\end{itemize}

它的安全性仅依赖于哈希函数，这使它成为后量子签名中的最后防线。至此，两大签名族——基于格的（ML-DSA、FN-DSA）和基于哈希的（SLH-DSA）——都已介绍完毕。下一部分将转向第三个相互独立的基础：基于编码的密码。

\printbibliography[heading=chapterrefs]
